% !TeX program = xelatex
% UTF-8 transcription of six handwritten EBV pages, in source page order.
% Original wording, abbreviations, numbers and highlights are retained.
% Source images: images/pic01.jpg through images/pic06.jpg.
% Compile directly with XeLaTeX, or input after loading xcolor and xepersian.
\ifdefined\settextfont
  \def\EBVStandalone{0}
\else
  \def\EBVStandalone{1}
  \documentclass[12pt,a4paper]{article}
  \usepackage[margin=2cm]{geometry}
  \usepackage{amsmath,amssymb,xcolor,enumitem}
  \usepackage{xepersian}
  \IfFontExistsTF{Amiri}{\settextfont{Amiri}}{\settextfont{FreeSerif}}
  \setlatintextfont{FreeSerif}
\fi
\providecommand{\en}[1]{\lr{#1}}
\newcommand{\ebvhl}[1]{\colorbox{yellow!35}{\strut #1}}
\newcommand{\ebvpage}[2]{\clearpage\section*{\en{EBV} ـ صفحهٔ #1}%
  {\footnotesize\en{#2}\par}}
\ifnum\EBVStandalone=1 \begin{document}\fi

\ebvpage{۱}{pic01.jpg — 20261011\_184339529.jpg}
\begin{itemize}
\item مونونوکلئوز عفونی: تب، گلودرد، لنفادنوپاتی، لنف آتیپیک (\ebvhl{لنفوسیتوز آتیپیک}).
  \begin{itemize}
  \item معمولاً در \ebvhl{بالغین جوان}.
  \item راه انتقال: ترشحات دهان، انتقال خون، پیوند \en{BM}.
  \item عفونت \en{EBV} در شیرخواران و کودکان خفیف یا بدون است.
  % The handwritten sentence ends with «خفیف یا بدون است».
  \end{itemize}
\item دورهٔ کمون: \en{4--6} هفته.
\item علائم پرودرومال: \en{1--2} هفته / خستگی، میالژی.
\item تب معمولاً \en{low-grade} و در هفتهٔ اول بیماری شایع است.
  \begin{itemize}
  \item حتی تا یک ماه هم ممکن است.
  \end{itemize}
\item لنفادنوپاتی: طی ۲ هفتهٔ اول بیماری بارزتر است.
  \begin{itemize}
  \item \ebvhl{معمولاً زنجیرهٔ خلفی گردن} / \ebvhl{می‌تواند ژنرالیزه}.
  \item معمولاً تندر، قرینه، \en{mobile}.
  \end{itemize}
\item فارنژیت: \ebvhl{بارزترین علامت} / بزرگی لوزه + اگزودا (مثل استرپتوکوک).
\item \ebvhl{اسپلنومگالی}: معمولاً در \ebvhl{هفتهٔ ۲ و ۳} بیماری.
\item تظاهرات پوستی:
  \begin{itemize}
  \item \ebvhl{راش پاپولر یا موربیلی‌فرم}: در ۵\% بیماران، بازو و تنه.
  \item \ebvhl{راش ماکولر}: اگر درمان با پنی‌سیلین.
  \item \ebvhl{اریتم ندوزوم}.
  \item \ebvhl{اریتم مولتی‌فرم}.
  \end{itemize}
\item طول مدت علائم \en{2--4} هفته است اما در ۱۰\% بیماران خستگی تا ۶ ماه می‌ماند.
\end{itemize}
% Printed notebook footer: GNOTE; handwritten page number: 1.

\ebvpage{۲}{pic02.jpg — 20261011\_184347491.jpg}
\begin{itemize}
\item شدت بیماری با تعداد سلول‌های \en{T} (\ebvhl{\en{CD8}}) و \ebvhl{\en{EBV DNA}} مرتبط است.
\item علائم \en{IMN} در افراد مسن:
  \begin{itemize}
  \item تب طول‌کشیده، خستگی، میالژی.
  \item فارنژیت، لنفادنوپاتی، اسپلنومگالی، لنفوسیتوز آتیپیک نادر است.
  \end{itemize}
\end{itemize}
\subsection*{بیماری لنفوپرولیفراتیو مرتبط با \en{EBV}}
\begin{itemize}
\item در افراد \ebvhl{نقص ایمنی مادرزادی} مثل نقص ایمنی مرکب شدید.
\item در افراد نقص ایمنی اکتسابی:
  \begin{itemize}
  \item \ebvhl{ایدز}.
  \item پیوند عضو یا مغز استخوان.
  \item داروی سرکوب ایمنی: لنفوم‌ها (\ebvhl{سلول \en{B}}).
  \end{itemize}
\item علائم: \ebvhl{تب}، \ebvhl{لنفادنوپاتی}، \ebvhl{علائم گوارشی}.
\item \ebvhl{سلول \en{B} آلوده به \en{EBV}} تکثیر یافته و وارد ارگان‌های مختلف می‌شود.
\end{itemize}
\subsection*{بیماری لنفوپرولیفراتیو وابسته به \en{X}}
\begin{itemize}
\item وابسته به \en{X} مغلوب، در پسران جوان، پس از عفونت با \en{EBV}.
\item بیماری لنفوپرولیفراتیو \ebvhl{کشنده}.
\item مرگ به علت مونونوکلئوز عفونی حاد.
\end{itemize}
\subsection*{آزمایشگاه}
\begin{itemize}
\item لنفوسیتوز در هفتهٔ ۲ و ۳: \ebvhl{\en{10,000--20,000}}.
  \begin{itemize}
  \item \ebvhl{\en{$>10\%$} لنف آتیپیک} (معمولاً لنف نوع \ebvhl{\en{CD8}}).
  \end{itemize}
\item \ebvhl{نوتروپنی خفیف} و \ebvhl{ترومبوسایتوپنی}: در ماه اول بیماری شایع است.
\item افزایش خفیف \ebvhl{\en{AST, ALT, ALK-P}}: در ۹۰\% بیماران.
\item افزایش بیلی‌روبین سرم: در ۴۰\% بیماران.
\end{itemize}
% Printed notebook footer: GNOTE; handwritten page number: 2.

\ebvpage{۳}{pic03.jpg — 20261011\_184355543.jpg}
\subsection*{تست آنتی‌بادی هتروفیل}
\begin{itemize}
\item در بیمار با \ebvhl{علائم \en{IMN}}، \ebvhl{لنف آتیپیک}، \ebvhl{تیتر \en{$\geq40$}} = عفونت \en{EBV}.
\item در هفتهٔ اول در ۴۰\% بیماران مثبت می‌شود.
\item در هفتهٔ سوم در \en{80--90\%} بیماران مثبت می‌شود.
\item انجام آزمایش‌های مکرر ضرورت دارد.
\item \en{Ab} می‌تواند \ebvhl{تا یک سال} مثبت بماند.
\item مواردی که \en{Ab} غیر قابل تشخیص است:
  \begin{itemize}
  \item کودکان زیر ۵ سال، افراد مسن، بیماران با علائم آتیپیک.
  \end{itemize}
\item نوع \ebvhl{\en{Monospot}} از نوع کلاسیک حساس‌تر است.
  \begin{itemize}
  \item حساسیت ۷۵\%، اختصاصیت ۹۰\%.
  \end{itemize}
\item مثبت کاذب تست \en{Monospot}:
  \begin{itemize}
  \item بیماری بافت همبند، لنفوم، هپاتیت ویروسی، مالاریا.
  \end{itemize}
\end{itemize}
\subsection*{تست آنتی‌بادی اختصاصی \en{EBV}}
\begin{itemize}
\item کاربرد:
  \begin{itemize}
  \item \ebvhl{بیمار مشکوک به \en{EBV} که فاقد \en{Ab} هتروفیل} است.
  \item مبتلایان به \ebvhl{عفونت آتیپیک}.
  \end{itemize}
\item آنتی‌بادی \en{IgM} و \en{IgG} علیه آنتی‌ژن کپسید (\ebvhl{\en{VCA}}).
  \begin{itemize}
  \item \en{IgM anti-VCA}: تشخیص \en{IMN} \ebvhl{حاد}.
  \item در \en{2--3} ماه اول بیماری افزایش می‌یابد.
  \item \en{IgG anti-VCA}: \ebvhl{تمام عمر} مثبت باقی می‌ماند.
  \end{itemize}
\item آنتی‌بادی \ebvhl{\en{EA-D}}:
  \begin{itemize}
  \item در ۷۰\% بیماران تولید می‌شود.
  \item معمولاً در \ebvhl{موارد شدید}.
  \item سطح آن با \ebvhl{کارسینوم نازوفارنژیال} ارتباط دارد.
  \item در عفونت مزمن بالاست.
  \item \en{3--9} ماه باقی می‌ماند.
  \end{itemize}
\end{itemize}
% Printed notebook footer: GNOTE; handwritten page number: 3.

\ebvpage{۴}{pic04.jpg — 20261011\_184404494.jpg}
\begin{itemize}
\item آنتی \en{EA-R}:
  \begin{itemize}
  \item در مبتلایان به \ebvhl{لنفوم بورکیت آفریقایی}.
  \item در عفونت مزمن فعال.
  \item جهت تشخیص مفید نیست.
  \end{itemize}
\item آنتی \en{IgA}: \ebvhl{جهت شناسایی بیماران کارسینوم نازوفارنژیال} و افراد در معرض خطر بالا برای این بیماری.
\item کشت خون، گلو، برای تشخیص عفونت حاد کاربردی ندارد چرا که \en{EBV} برای تمام عمر در اوروفارنکس و سلول‌های \en{B} باقی می‌ماند.
\item تغییرات آزمایشگاهی در \en{IMN}:
  \begin{itemize}
  \item لکوسیتوز به صورت لنفوسیتوز.
  \item نوتروپنی خفیف.
  \item ترومبوسایتوپنی.
  \item \ebvhl{آنمی همولیتیک اتوایمیون} (\en{AIHA}).
  \end{itemize}
\end{itemize}
\subsection*{\en{DDX}}
\begin{itemize}
\item داروهایی که می‌توانند در \en{DDX} با \en{IMN} قرار بگیرند:
  \begin{itemize}
  \item فنی‌توئین، کاربامازپین، سولفونامیدها، مینوسایکلین.
  \end{itemize}
\item شایع‌ترین علت \en{IMN} هتروفیل منفی، \ebvhl{\en{CMV}} است.
\end{itemize}
\subsection*{درمان \en{IMN}}
\begin{itemize}
\item حمایتی: استراحت، مسکن، عدم فعالیت شدید بدنی در ماه اول (پارگی طحال).
\item \en{CS} خوراکی (\ebvhl{پردنیزولون}): \en{40--60 mg} در روز برای \en{3--4} روز، با کاهش دوز طی \en{1--2} هفته.
  \begin{itemize}
  \item اگر \ebvhl{هایپرتروفی لوزه شدید}، \ebvhl{\en{AIHA}}، لنفوهیستیوسیتوز هموفاگوسیتیک، \ebvhl{ترومبوسایتوپنی شدید}، تب و کسالت شدید، \ebvhl{درگیری \en{CNS}}، \ebvhl{درگیری قلبی}.
  \end{itemize}
\end{itemize}
% Printed notebook footer: GNOTE; handwritten page number: 4.

\ebvpage{۵}{pic05.jpg — 20261011\_184413363.jpg}
\begin{itemize}
\item آسیکلوویر در \en{EBV} اثر چندانی ندارد.
\item آسیکلوویر در درمان \en{IMN} مؤثر نیست ولی در درمان \ebvhl{لوکوپلاکی مویی دهانی} مؤثر است.
\end{itemize}
\subsection*{بیماری لنفوپرولیفراتیو \en{EBV} پس از پیوند}
\begin{itemize}
\item معمولاً به درمان ضد ویروسی جواب نمی‌دهد.
\item درمان:
  \begin{itemize}
  \item \ebvhl{کاهش سرکوب ایمنی}.
  \item \ebvhl{\en{anti-CD20 Ab}}: \ebvhl{ریتوکسیمب}.
  \end{itemize}
\end{itemize}
\subsection*{عوارض}
\begin{itemize}
\item مرگ در اثر \en{IMN} نادر است و معمولاً بر اثر درگیری \en{CNS}، پارگی طحال، انسداد راه هوایی و عفونت باکتریال ثانویه.
\item عوارض \en{CNS}: معمولاً طی ۲ هفتهٔ اول عفونت.
  \begin{itemize}
  \item \ebvhl{مننژیت و انسفالیت}: سردرد، مننژیسموس، آتاکسی مخچه‌ای.
  \item در موارد عفونت حاد: فلج اعصاب کرانیال (شایع \ebvhl{\en{VII}})، سندرم گیلن‌باره، میلیت عرضی حاد، نوریت محیطی.
  \end{itemize}
\item عوارض خونی:
  \begin{itemize}
  \item در ۲\% موارد \ebvhl{\en{AIHA}} (طی ۲ هفتهٔ اول).
  \item در بسیاری موارد \ebvhl{کومبس مثبت} است و \ebvhl{آگلوتینین سرد} علیه \ebvhl{آنتی‌ژن \en{RBC}} وجود دارد.
  \end{itemize}
\item پارگی شکم:
  % Source heading reads «پارگی شکم»; preserved without medical correction.
  \begin{itemize}
  \item در ۵\% موارد، در مردان شایع‌تر.
  \item درد شکم، درد ارجاعی شانه، هیپوتانسیون (\ebvhl{\en{BP$\downarrow$}}).
  \end{itemize}
\item هایپرتروفی لنفوئید در لوزه‌ها و آدنوئید:
  \begin{itemize}
  \item انسداد راه هوایی فوقانی.
  \end{itemize}
\end{itemize}
% Printed notebook footer: GNOTE; handwritten page number: 5.

\ebvpage{۶}{pic06.jpg — 20261011\_184424912.jpg}
\begin{itemize}
\item \en{EBV} ریسک برخی بدخیمی‌ها را بالا می‌برد:
  \begin{itemize}
  \item کارسینوم نازوفارنژیال، کارسینوم معده، لنفوم بورکیت، لنفوم هوچکین، لنفوم \en{T}.
  \item \ebvhl{لنفوم \en{B}} (\ebvhl{در مبتلایان به نقص ایمنی}).
  \item \ebvhl{تومور عضلهٔ صاف} (\ebvhl{در مبتلایان به نقص ایمنی}).
  \end{itemize}
\end{itemize}
% Printed notebook footer: GNOTE; handwritten page number: 6.

\ifnum\EBVStandalone=1
  \let\EBVEnd\enddocument
\else
  \let\EBVEnd\relax
\fi
\EBVEnd
